\documentclass[UTF8,a4paper,11pt,openany]{ctexbook}
\usepackage{../common/statlearnbook}
\SLSetBookMetadata{第 5 册}{统计学习方法讲义 v2.7.0}{采样方法、主题模型与图排序}
\begin{document}
\setcounter{page}{595}
\setcounter{chapter}{32}
\setcounter{figure}{0}
\pagestyle{headings}
\chaptermark{马尔可夫链蒙特卡罗法}
\sectionmark{本章局部知识依赖}
图\ref{fig:V5-C03-dependency}要求分支阅读：先沿“MH构造$\to K\to$细致平衡证书”核对
$\pi(\mathrm dx)K(x,\mathrm dy)=\pi(\mathrm dy)K(y,\mathrm dx)$只推出$\pi K=\pi$；
再分别沿链结构分支判断遍历时间平均与边缘分布收敛所需的附加条件，并沿诊断分支检查轨迹相关性。
% v2.7.0 figure UID: FIG-P596-01
% R2 semantic repair: separate stationarity, ergodic averages, marginal convergence, and diagnostics.
\pgfkeysifdefined{/tikz/alt}{}{\tikzset{alt/.code={}}}
\tikzset{slfig-FIG-P596-01/.style={font=\fontsize{9.6pt}{11.5pt}\selectfont,every path/.append style={line cap=round,line join=round}}}
\begin{figure}[htbp]
\centering
\begin{tikzpicture}[
  alt={对象—关系—结论：目标分布$\pi$与提议$q$经接受率$\alpha$构造MH转移核$K$；若$K$满足细致平衡$\pi(\mathrm dx)K(x,\mathrm dy)=\pi(\mathrm dy)K(y,\mathrm dx)$，则只推出$\pi K=\pi$，即$\pi$平稳。遍历时间平均还须另行核验有限状态不可约或一般状态正Harris常返；边缘分布收敛还须在相应结构条件上加非周期性。$K$通常形成相关轨迹，但独立核$K(x,\mathrm dy)=\pi(\mathrm dy)$是反例，相关性应由轨迹、自相关、MCSE与ESS诊断。},
  slfig-FIG-P596-01,>=Stealth,
  every node/.style={font=\fontsize{9.6pt}{11.5pt}\selectfont,align=center},
  nodebox/.style={draw=SLRuleGray,fill=white,inner xsep=3.5pt,inner ysep=3.0pt,minimum height=8.5mm,line width=.72pt},
  inputbox/.style={nodebox,rounded corners=5pt,draw=SLGold,fill=SLGold!7,line width=.90pt},
  constructbox/.style={nodebox,rounded corners=5pt,draw=SLGold,fill=SLGold!11,line width=.90pt},
  kernelbox/.style={circle,draw=SLBlue,fill=SLBlue!8,minimum size=14mm,inner sep=1.5pt,line width=.95pt},
  certificatebox/.style={nodebox,rounded corners=0pt,draw=SLTeal,fill=SLTeal!6,line width=.82pt,densely dashed},
  resultbox/.style={nodebox,rounded corners=0pt,draw=SLBlue,fill=SLBlue!7,line width=.95pt},
  conditionbox/.style={nodebox,rounded corners=0pt,draw=SLBlue,fill=SLBlue!4,line width=.82pt,dash pattern=on 4pt off 1.5pt on .9pt off 1.5pt},
  diagnosticbox/.style={nodebox,rounded corners=5pt,draw=SLGray,fill=SLSoftGray,line width=.82pt,densely dotted},
  scopebox/.style={nodebox,rounded corners=2pt,draw=SLRuleGray,fill=SLSoftGray,line width=.68pt},
  construction/.style={draw=SLGold,line width=.95pt},
  certificate/.style={draw=SLTeal,line width=.82pt,densely dashed},
  condition/.style={draw=SLBlue,line width=.82pt,dash pattern=on 4pt off 1.5pt on .9pt off 1.5pt},
  diagnosis/.style={draw=SLGray,line width=.82pt,densely dotted}
]
  % MH construction and its stationarity certificate.
  \node[inputbox,text width=19mm] (input) at (-5.70,2.75) {目标 $\pi$\\提议 $q$};
  \node[constructbox,text width=20mm] (mh) at (-2.90,2.75) {接受率 $\alpha$\\MH 构造};
  \node[kernelbox] (kernel) at (-.55,2.75) {核 $K$};
  \node[certificatebox,text width=29mm] (balance) at (2.20,2.75)
    {细致平衡证书\\$\pi(\mathrm dx)K(x,\mathrm dy)$\\$=\pi(\mathrm dy)K(y,\mathrm dx)$};
  \node[resultbox,text width=20mm] (stationary) at (5.55,2.75) {$\pi K=\pi$\\$\pi$ 平稳};
  \draw[->,construction] (input.east)--(mh.west);
  \draw[->,construction] (mh.east)--(kernel.west);
  \draw[->,certificate] (kernel.east)--(balance.west);
  \draw[->,certificate] (balance.east)--(stationary.west);

  \node[scopebox,text width=28mm] at (6.00,1.55)
    {细致平衡止于平稳\\不证不可约 / 遍历};

  % Separate branches for diagnostics and for convergence statements.
  \node[diagnosticbox,text width=37mm] (trajectory) at (-3.65,.35)
    {通常形成相关轨迹\\独立核 $K(x,\mathrm dy)$\\$=\pi(\mathrm dy)$ 是反例};
  \node[diagnosticbox,text width=35mm] (diagnostics) at (-3.65,-1.50)
    {相关性需诊断\\轨迹 / 自相关 / MCSE / ESS};
  \node[conditionbox,text width=34mm] (structure) at (2.30,.35)
    {另行核验链结构\\有限状态：不可约\\一般状态：正 Harris 常返};
  \node[resultbox,text width=20mm] (average) at (5.70,.35) {遍历\\时间平均};
  \node[conditionbox,text width=31mm] (margcond) at (2.30,-1.50)
    {相应结构条件\\再加非周期性};
  \node[resultbox,text width=20mm] (marginal) at (5.70,-1.50) {边缘分布\\收敛};

  \draw[->,diagnosis] (kernel.south west) |- (trajectory.east);
  \draw[->,diagnosis] (trajectory.south)--(diagnostics.north);
  \draw[->,condition] (kernel.south east) |- (structure.west);
  \draw[->,condition] (structure.east)--(average.west);
  \draw[->,condition] (structure.south)--(margcond.north);
  \draw[->,condition] (margcond.east)--(marginal.west);

  % Grayscale-readable line samples; color is only a redundant cue.
  \draw[construction] (-5.95,-2.85)--(-5.05,-2.85);
  \node[anchor=west] at (-4.85,-2.87) {实线：MH 构造};
  \draw[certificate] (.20,-2.85)--(1.10,-2.85);
  \node[anchor=west] at (1.30,-2.87) {虚线：平稳证书};
  \draw[condition] (-5.95,-3.50)--(-5.05,-3.50);
  \node[anchor=west] at (-4.85,-3.52) {点划线：结构条件};
  \draw[diagnosis] (.20,-3.50)--(1.10,-3.50);
  \node[anchor=west] at (1.30,-3.52) {点线：相关性诊断};
\end{tikzpicture}
\caption{细致平衡只验证平稳性；遍历平均和分布收敛还需相应链结构条件。}\label{fig:V5-C03-dependency}
\end{figure}

\FloatBarrier
\noindent\textbf{读图检查（仅针对图\ref{fig:V5-C03-dependency}）。}只沿细致平衡分支能得到$\pi K=\pi$，不能得到不可约或遍历；时间平均须另行核验有限状态不可约或一般状态正Harris常返，边缘分布收敛还须在相应结构条件上核验非周期性。图中的“通常相关轨迹”只是诊断动机；独立核$K(x,\mathrm dy)=\pi(\mathrm dy)$说明相关性不是任意转移核的必然结论。
\end{document}
